% Length scales and friction parameters shared by the cohesive-zone,
% linear-slip-weakening and rate-and-state descriptions of the PMMA fault.
% Values are produced by CohesiveZoneModel/Lc_estimate.py.
\documentclass[11pt]{article}

\usepackage[a4paper, margin=1.3cm]{geometry}
\usepackage{booktabs}
\usepackage{array}
\usepackage{amsmath}
\usepackage{siunitx}
\usepackage[hidelinks]{hyperref}

% Ragged-right paragraph column, so the fixed-width cells below do not stretch.
\newcolumntype{L}[1]{>{\raggedright\arraybackslash}p{#1}}
\newcommand{\doi}[1]{\href{https://doi.org/#1}{\scriptsize\texttt{#1}}}
\title{From cohesive zone to slip weakening and rate and state}

\begin{document}
\maketitle

The cohesive-zone model is the anchor. It fixes the process-zone size $X_c$; the fracture energy $\Gamma$, the linear-slip-weakening (LSW) parameters and the rate-and-state (RSF) parameters all follow from it at the current friction coefficients. $D_c$ is therefore a derived quantity, and carrying an old $D_c$ across a change of $\mu_k$ silently rescales $\Gamma$.

Two kinds of quantity appear below and must not be confused. A \emph{slip} is a relative displacement across the fault, measured along the sliding direction; a \emph{fault length} is a distance measured along the fault plane. Both carry units of length, but they differ here by roughly four orders of magnitude. Conventions differ between sources as well: $L_c$ and $h^{*}$ are half-lengths, whereas $X_c$ and $L_b$ are full zone lengths. Comparing a half-length against a physical patch width requires the factor of two.

\begin{table}[htbp]
  \centering
  \footnotesize
  \caption{Length scales and friction parameters of the PMMA fault, anchored on
  $X_c = \SI{5}{\milli\meter}$ with $\mu_s = 0.8$, $\mu_k = 0.45$,
  $\sigma_n = \SI{16}{\mega\pascal}$, $E = \SI{7.662}{\giga\pascal}$ and
  $\nu = 0.2$. RSF values are quoted for the middle (rupture-carrying) zone,
  $a = 0.005$, $b = 0.02581940$.}
  \label{tab:czm-lsw-rsf}
  \setlength{\tabcolsep}{4pt}
  \begin{tabular}{@{}l L{2.6cm} l l l r L{2.5cm} L{4.0cm}@{}}
    \toprule
    Symbol & Quantity & Model & Type & Conv. & Value & Source & DOI \\
    \midrule
    $X_c$
      & Cohesive zone size
      & CZM & Length & Full
      & \SI{5}{\milli\meter}
      & Palmer \& Rice (1973); Kammer \& McLaskey (2019)
      & \doi{10.1098/rspa.1973.0040} \newline
        \doi{10.1016/j.epsl.2019.01.031} \\
    \addlinespace
    $\Gamma$
      & Fracture energy
      & Shared & Energy & ---
      & \SI{22.2348}{\joule\per\meter\squared}
      & $\Gamma = \Delta\tau D_c / 2$
      & --- \\
    \addlinespace
    $D_c$
      & Characteristic slip
      & LSW & \textbf{Slip} & ---
      & \SI{7.94099e-3}{\milli\meter}
      & Inverted from $X_c$
      & \doi{10.1098/rspa.1973.0040} \\
    \addlinespace
    $L_c$
      & Critical nucleation length
      & LSW & Length & \textbf{Half}
      & \SI{3.60253}{\milli\meter}
      & Andrews (1976)
      & \doi{10.1029/JB081i032p05679} \\
    $2L_c$
      & Smallest unstable patch
      & LSW & Length & Full
      & \SI{7.20506}{\milli\meter}
      & Andrews (1976)
      & \doi{10.1029/JB081i032p05679} \\
    \addlinespace
    $D_c^{\mathrm{RSF}}$
      & State evolution distance
      & RSF & \textbf{Slip} & ---
      & \SI{3.76505e-4}{\milli\meter}
      & Ageing law, matched on $\Gamma$
      & \doi{10.1029/JB084iB05p02161} \\
    \addlinespace
    $L_b$
      & Process zone size
      & RSF & Length & Full
      & \SI{3.63702}{\milli\meter}
      & Rubin \& Ampuero (2005)
      & \doi{10.1029/2005JB003686} \\
    \addlinespace
    $h^{*}_{\mathrm{RR}}$
      & Nucleation length, linear stability
      & RSF & Length & \textbf{Half}
      & \SI{3.54253}{\milli\meter}
      & Rice \& Ruina (1983)
      & \doi{10.1115/1.3167042} \\
    $h^{*}_{\mathrm{RA}}$
      & Nucleation length, ageing law
      & RSF & Length & \textbf{Half}
      & \SI{3.56108}{\milli\meter}
      & Rubin \& Ampuero (2005)
      & \doi{10.1029/2005JB003686} \\
    \bottomrule
  \end{tabular}
\end{table}


\newpage
\section*{Governing expressions}

With $\Delta\tau = (\mu_s - \mu_k)\sigma_n$ and the plane-strain combination
$\mu^{*} = \mu/(1-\nu)$, the entries above follow from
%
\begin{align}
  X_c &= \frac{9\pi}{16}\,\frac{\mu^{*}\Gamma}{\Delta\tau^{2}},
  &
  \Gamma &= \tfrac{1}{2}\Delta\tau D_c,
  \\
  L_c &= \frac{4}{\pi}\,\frac{\mu^{*}\Gamma}{\Delta\tau^{2}},
  &
  L_b &= \frac{\mu^{*} D_c^{\mathrm{RSF}}}{b\,\sigma_n},
  \\
  h^{*}_{\mathrm{RR}} &=
    \frac{\pi}{4}\,\frac{\mu^{*} D_c^{\mathrm{RSF}}}{(b-a)\,\sigma_n},
  &
  h^{*}_{\mathrm{RA}} &=
    \frac{2}{\pi}\,\frac{\mu^{*}\, b\, D_c^{\mathrm{RSF}}}{(b-a)^{2}\sigma_n}.
\end{align}
%
$h^{*}_{\mathrm{RR}}$ applies for $a/b$ well below $0.378$ and
$h^{*}_{\mathrm{RA}}$ as $a/b$ approaches unity; both are undefined where the
zone is velocity neutral or strengthening. The two nucleation estimates and
$L_c$ answer different questions --- a linear-stability size for a uniform RSF
fault against an energy balance for a slip-weakening crack --- so their close
agreement here should be read as coincidence rather than as mutual
confirmation.

\section*{Rate-and-state coefficients}

The RSF friction coefficient is
$\mu = f_0 + a\ln(V/V^{*}) + b\ln(V^{*}\theta/D_c^{\mathrm{RSF}})$, so $a$ is
the instantaneous (direct) response to a change of slip rate and $b$ is the
evolution effect that follows as the state variable $\theta$ relaxes over a
slip of order $D_c^{\mathrm{RSF}}$. The steady-state combination $a-b$ sets
stability.

\begin{table}[htbp]
  \centering
  \footnotesize
  \caption{Zone-wise rate-and-state coefficients. The friction drop alone fixes
  $a-b = (\mu_k-\mu_s)/\ln(V_2/V_1) = -0.0208194$ for the calibrated zone; $a$
  only decides how that drop is split between the direct and evolution terms.
  Velocities are $V_1 = \SI{1e-4}{\milli\meter\per\second}$ and
  $V_2 = \SI{2000}{\milli\meter\per\second}$.}
  \label{tab:rsf-zones}
  \begin{tabular}{@{}l r r r r l@{}}
    \toprule
    Zone & $a$ & $b$ & $a-b$ & $\mu(V_2)$ & Behaviour \\
    \midrule
    Loading & 0.004 & 0.004        & $0$          & 0.8000 & Velocity neutral \\
    Middle  & 0.005 & 0.02581940   & $-0.0208194$ & 0.4500 & Velocity weakening \\
    Leading & 0.010 & 0.004        & $+0.0060000$ & 0.9009 & Velocity strengthening \\
    \bottomrule
  \end{tabular}
\end{table}

Only the middle zone carries the rupture, so it is the one calibrated to reach
$\mu_k$ at $V_2$, and the shared $D_c^{\mathrm{RSF}}$ is matched on it. The
loading zone transmits load without weakening, and the leading zone acts as a
barrier. Because $L_b \propto b^{-2}$, the middle zone also binds the mesh: at
$L_b = \SI{3.63702}{\milli\meter}$ a \SI{0.5}{\milli\meter} cell gives $7.3$
cells across the process zone, against a usual floor of five. $L_b$ is the
quasi-static width and contracts as the rupture accelerates, so that margin
should not be spent.

\newpage
\section*{Why the two characteristic slips differ by an order of magnitude}

Table~\ref{tab:czm-lsw-rsf} carries two slips that share a fracture energy yet
differ by a factor of $21$: $D_c = \SI{7.94e-3}{\milli\meter}$ for linear slip
weakening and $D_c^{\mathrm{RSF}} = \SI{3.77e-4}{\milli\meter}$ for rate and
state. They are not two estimates of one quantity. They measure different
things.

\paragraph{What each slip means.}
In linear slip weakening, $D_c$ is the slip at which weakening \emph{ends}: the
strength falls linearly from $\tau_p$ to $\tau_r$ and is then flat, so the whole
breakdown work $\Gamma = \tfrac12\Delta\tau D_c$ is spent by the time the fault
has slipped $D_c$. In rate and state, the slip written $D_c^{\mathrm{RSF}}$ (the
$L$ of Dieterich and Ruina) is the distance over which the state variable
$\theta$ is renewed. Under the ageing law $\dot\theta = 1 - V\theta/L$, a step
in slip rate leaves $\theta$ relaxing as $\exp(-\delta/L)$; $L$ is the
$e$-folding slip of the contact population, on the order of an asperity
diameter. It marks where weakening \emph{begins} to be felt, not where it is
complete, because the friction depends on $\theta$ only through a logarithm.

\paragraph{Where the factor comes from.}
Start every point at the steady state of the initial rate $V_1$, so
$\theta_0 = L/V_1$, and step the rate to $V_2$. With $W = \theta_0 V_2/L = V_2/V_1$
the strength excess above the new steady value is
%
\begin{equation}
  \tau(\delta) - \tau_{\mathrm{ss}}
    = \sigma_n\, b \ln\!\bigl[1 + (W-1)\,e^{-\delta/L}\bigr],
  \label{eq:excess}
\end{equation}
%
which starts at $\sigma_n b\ln W$ and decays only logarithmically. Its integral
over slip is the breakdown work,
%
\begin{equation}
  \Gamma = \sigma_n\, b\, L\, K(W),
  \qquad
  K(W) = -\mathrm{Li}_2\bigl(-(W-1)\bigr)
       \simeq \tfrac12 \ln^{2} W + \tfrac{\pi^{2}}{6},
  \label{eq:gamma-rsf}
\end{equation}
%
and $K = 142.95$ at the ratio $W = 2\times10^{7}$ the cases use. If the same
curve is read as slip weakening with the same peak-to-residual drop
$\Delta\tau_{\mathrm{RSF}} = \sigma_n b \ln W$, the equivalent LSW slip is
$D_c^{\mathrm{eq}} = 2\Gamma/\Delta\tau_{\mathrm{RSF}}$, so
%
\begin{equation}
  \frac{D_c^{\mathrm{eq}}}{L}
    = \frac{2K(W)}{\ln W}
    \simeq \ln W + \frac{\pi^{2}}{3\ln W}
    = 17.0 .
  \label{eq:ratio}
\end{equation}
%
The factor is set by the logarithm of the velocity jump alone. It is the
$D_0^{\mathrm{eq}}/L \sim 15$ that Cocco and Bizzarri obtained numerically for
crustal parameters; the residual difference is their slightly smaller $\ln W$
and their reading of $D_c^{\mathrm{eq}}$ off a computed traction curve rather
than from the energy.

The table's ratio of $21.1$ carries one more factor. It compares the
rate-and-state $L$ against the LSW $D_c$ inverted from the cohesive-zone
$\Delta\tau = (\mu_s - \mu_k)\sigma_n = \SI{5.60}{\mega\pascal}$, whereas the
rate-and-state curve overshoots through the direct effect to a peak
$\mu = \mu_s + a\ln W = 0.884$ and drops
$\Delta\tau_{\mathrm{RSF}} = \SI{6.94}{\mega\pascal}$. Hence
$21.1 = 17.0 \times 6.94/5.60$. Both numbers are correct; they answer different
questions.

\begin{table}[htbp]
  \centering
  \footnotesize
  \caption{Fraction of the rate-and-state breakdown work spent within a given
  multiple of $L$, from Eq.~\eqref{eq:excess} at $W = 2\times10^{7}$. Only a
  ninth of $\Gamma$ is dissipated over the first $L$; the linear-slip-weakening
  picture, in which all of it is spent by $D_c$, corresponds to the whole tail.}
  \label{tab:gamma-partition}
  \begin{tabular}{@{}r r@{}}
    \toprule
    Slip & Share of $\Gamma$ \\
    \midrule
    $1\,L$    & \SI{11.4}{\percent} \\
    $2\,L$    & \SI{22.1}{\percent} \\
    $5\,L$    & \SI{50.1}{\percent} \\
    $10\,L$   & \SI{82.6}{\percent} \\
    $11.8\,L$ & \SI{90.0}{\percent} \\
    $17\,L$   & \SI{99.5}{\percent} \\
    \bottomrule
  \end{tabular}
\end{table}

\paragraph{Consequence for the model.}
The slip that governs the mesh and the nucleation size is $L$, not
$D_c^{\mathrm{eq}}$: $L_b$ and $h^{*}$ in Table~\ref{tab:czm-lsw-rsf} are built
on $L$, in line with the finding of Cocco and Bizzarri that the nucleation patch
scales with $L$. Matching $\Gamma$ to the cohesive zone therefore forces
$L$ to be small, and that in turn is what sets the \SI{0.5}{\milli\meter} cell
size. The literature on crustal faults reaches the same conclusion from the
other side: the $D_c$ recovered from dynamic or kinematic models is an outcome
of the rupture and of the modelling assumptions, not a constitutive constant,
and its apparent scaling with final slip can be an artefact of poorly
constrained kinematic parameters.

\begin{table}[htbp]
  \centering
  \footnotesize
  \caption{What the rate-and-state versus slip-weakening literature establishes
  about the two characteristic slips.}
  \label{tab:dc-literature}
  \setlength{\tabcolsep}{4pt}
  \begin{tabular}{@{}L{3.2cm} L{8.9cm} L{4.9cm}@{}}
    \toprule
    Reference & Finding & DOI \\
    \midrule
    Cocco \& Bizzarri (2002), \emph{GRL} 29(11), 1516
      & A rate-and-state fault produces an effective slip-weakening traction
        curve whose equivalent distance $D_0^{\mathrm{eq}}$ differs from the
        state-evolution slip $L$, with $D_0^{\mathrm{eq}}/L \sim 15$. The
        nucleation patch scales with $L$, not with $D_0^{\mathrm{eq}}$.
      & \doi{10.1029/2001GL013999} \\
    \addlinespace
    Bizzarri \& Cocco (2003), \emph{JGR} 108(B8), 2373
      & 2-D in-plane spontaneous ruptures under rate and state: the equivalent
        slip-weakening distance is a function of $L$ \emph{and} of the initial
        mechanical state, so it is an outcome of the rupture rather than a
        material constant.
      & \doi{10.1029/2002JB002198} \\
    \addlinespace
    Tinti et al.\ (2004), \emph{GRL} 31, L02611
      & Estimates of $D_c$ from the traction evolution of 2-D in-plane cracks
        under slip-weakening and rate-and-state laws are compared; the recovered
        $D_c$ depends on the constitutive law and on how the breakdown is read
        off the curve.
      & \doi{10.1029/2003GL018811} \\
    \addlinespace
    Tinti et al.\ (2009), \emph{GJI} 177, 1205
      & $D_c$ inferred from kinematic slip-velocity functions is highly
        sensitive to the assumed rise time and acceleration duration; imposing
        wrong but constant values manufactures an apparent $D_c$--final-slip
        correlation absent from the target models.
      & \doi{10.1111/j.1365-246X.2009.04143.x} \\
    \addlinespace
    Cocco, Tinti, Marone \& Piatanesi (2009), \emph{Int.\ Geophys.} 94, 163
      & Review of the scaling of $D_c$ with final slip and of the physical
        processes (thermal pressurisation, flash heating, gouge production)
        that make the breakdown distance scale dependent rather than a fixed
        fault property.
      & \doi{10.1016/S0074-6142(08)00007-7} \\
    \bottomrule
  \end{tabular}
\end{table}

\end{document}
